\section{Multi-Agent Subagent System}
\label{sec:subagents}

OpenSkynet includes a multi-agent framework for decomposing complex tasks across specialized subagents, each with restricted permissions and optionally isolated execution environments.

\subsection{Subagent Definition}

Subagents are defined as Markdown files with YAML frontmatter, making them easy to create, share, and customize. Each subagent specification includes: a name, description, mode (defaulting to ``subagent''), an optional model override, a permission map (tool $\rightarrow$ allow/ask/deny), a maximum iteration limit, and a system prompt body. OpenSkynet ships with 7 built-in subagents:

\begin{itemize}[leftmargin=*,itemsep=1pt]
    \item \textbf{Browser} (max 8 iterations): Web browsing, form filling, data extraction. Denied: terminal, file write.
    \item \textbf{Code} (max 30 iterations): Full-featured coding agent with a 5-stage workflow (Explore $\rightarrow$ Plan $\rightarrow$ Execute $\rightarrow$ Verify $\rightarrow$ Summarize). Denied: browser.
    \item \textbf{Debug} (max 5 iterations): Diagnostic agent for analyzing failures. Read-only access to files and terminal.
    \item \textbf{Explore} (max 3 iterations): Rapid survey agent. Max 5 bullet points in summary. No file writes.
    \item \textbf{Integrate} (max 50 iterations): Merge feature branches and resolve conflicts. Full terminal access.
    \item \textbf{Review} (max 3 iterations): Critique and verify existing work. Rates as PASS/NEEDS\_FIX/REJECT.
    \item \textbf{RedTeam} (max 10 iterations): Adversarial test engineer. Creates edge-case tests (null, empty, large, Unicode, concurrent).
\end{itemize}

User-defined subagents can be added to \texttt{\textasciitilde/.openskynet/agents/} and may override built-in agents by name.

\subsection{Permission-Based Isolation}

Each subagent inherits a permission map that filters the global tool registry before execution. Permissions are resolved in priority order: deny (by skill name) $\rightarrow$ deny (by source) $\rightarrow$ always\_allow\_source $\rightarrow$ always\_allow\_skill $\rightarrow$ allow\_once $\rightarrow$ skip $\rightarrow$ none. When a tool's permission is ``ask'', an async approval callback is invoked, enabling human-in-the-loop approval for high-risk actions.

This enforces the principle of least privilege: a browser subagent cannot execute terminal commands; a code subagent cannot navigate web pages; and a review agent cannot modify files. Risk assessment (\texttt{assess\_risk}) classifies tool calls as low, medium, or high risk based on command patterns (safe command whitelist, risky URL patterns, destructive tool identifiers), enabling additional gating for dangerous operations.

\subsection{Parallel Execution}

When the Manager Agent decomposes a task into multiple independent subtasks (Decompose strategy), the \texttt{SubagentFactory} spawns up to 3 subagents concurrently using \texttt{asyncio.Semaphore} for concurrency control. Each subagent can optionally receive an isolated browser context (temporary directory, headless, stealth mode) to prevent session interference. Results are gathered in order and merged.

A shared scratchpad---a versioned key-value store---enables subagents to communicate context across execution boundaries. Parent context (task description, recent errors, recent observations) is injected into each subagent's system prompt via \texttt{<parent\_task>}, \texttt{<parent\_errors>}, and \texttt{<parent\_observations>} tags.

\subsection{Nesting Control}

To prevent infinite recursion, the \texttt{SubagentFactory} tracks nesting depth and refuses to spawn beyond a configurable maximum depth. This is essential when a code subagent delegates a subtask to another code subagent, which could otherwise cascade indefinitely.
